\section{Related Work}
Indonesian slang normalization addresses the lexical variation found in informal text. The Colloquial Indonesian Lexicon provides informal-to-formal mappings, including entries with multiple alternatives~\cite{salsabila2018}. Such ambiguity matters because a dictionary substitution does not necessarily preserve the meaning intended in a particular review. IndoCollex evaluates normalization through annotated word transformations~\cite{wibowo2021}; success on that task does not by itself establish improved sentiment classification. Downstream evaluation must therefore assess both the benefits and errors associated with replacing informal forms.

The classifier's representation is another important factor. IndoNLU provides SmSA and the IndoBERT pre-trained models~\cite{wilie2020}. Separately, IndoBERTweet extends an Indonesian BERT model with a vocabulary adapted to Twitter~\cite{koto2021}. These resources highlight the need to examine normalization in conjunction with the model's handling of informal language. Siva et al.~\cite{siva2026} compare IndoBERT with an SVM that uses word and character TF-IDF. However, this combined representation does not reveal the separate responses of word and character features to normalization. The use of word-level matching, character overlap, and contextual subword representations all lead to different hypotheses about the impact of substitution, each requiring direct comparison.

Recent work does provide direct evidence on the effect of preprocessing. Bustamin et al.~\cite{bustamin2025} report higher IndoBERT accuracy on tweets when dictionary normalization is combined with edit-distance correction (the combined approach). Budianoor et al.~\cite{budianoor2026} evaluate slang normalization against a minimal-cleaning baseline on culinary comments and see a macro-F1 go from 0.8272 to 0.8482. Prichanti et al.~\cite{prichanti2026} find different preferences for heavy versus minimal preprocessing for Word2Vec/logistic regression and IndoRoBERTa. Their caveat that they don't have a component-wise ablation means they can't really attribute the differences they see to slang replacement specifically. All of this serves to underline the need to control the intervention and evaluate it across representations.

Evaluation design also has an impact on interpretation. Babanejad et al.~\cite{babanejad2020} make a distinction between preprocessing as part of the representation learning and preprocessing for downstream tasks. Kuswanto and Tanamal~\cite{kuswanto2026} use multiple seeds and paired uncertainty estimates in their Indonesian app review comparison work and look at the effect of different preprocessing configurations for the emotion classification task. The use of repeated runs and the comparison of different preprocessing options is therefore well established. Bestgen~\cite{bestgen2022} stresses the importance of reporting differences with confidence intervals, which enables direct comparison of the effect of normalization rather than any conclusions being drawn from the separate model scores alone.

Together these studies lead to a controlled evaluation which disentangles representation, preprocessing policy and the point at which substitution occurs. The present study evaluates one fixed dictionary intervention, separate word, character and contextual models and shared audited splits on SmSA and PRDECT-ID~\cite{sutoyo2022}. This evaluation is possible using existing resources and consolidation efforts (e.g. NusaCrowd~\cite{cahyawijaya2023}) without any new scraping or annotation. The contribution of the work is a joint evaluation of the separate pipelines and input changes to the fixed models, with paired uncertainty and error transitions, in two defined sentiment regimes.
